\begin{table}[H]
\centering
\caption{Niveles de remuneración laboral por departamento, 2025}
\label{tab:dept_remuneracion_niveles_24}
\scriptsize
\begin{tabular}{lrrrrrr}
\toprule
Departamento & Rem./trab. & Rel. líder & Rem./hora & Rel. líder & Ocupados & Horas/trab. \\
\midrule
Bogotá D.C. & 2,89 & 100,0\% & 14,9 & 100,0\% & 4,16 & 44,7 \\
Antioquia & 2,18 & 75,4\% & 11,2 & 75,2\% & 3,31 & 44,8 \\
Cundinamarca & 1,99 & 68,6\% & 10,4 & 69,8\% & 1,71 & 44,0 \\
Santander & 1,94 & 66,9\% & 9,8 & 65,9\% & 1,05 & 45,4 \\
Valle del Cauca & 1,91 & 65,9\% & 10,0 & 67,0\% & 2,01 & 44,0 \\
Risaralda & 1,82 & 62,9\% & 9,4 & 62,6\% & 0,45 & 45,0 \\
Meta & 1,80 & 62,1\% & 8,9 & 59,4\% & 0,52 & 46,8 \\
Caldas & 1,78 & 61,7\% & 9,4 & 63,2\% & 0,47 & 43,6 \\
Quindío & 1,71 & 59,0\% & 9,2 & 61,6\% & 0,26 & 42,8 \\
Atlántico & 1,61 & 55,8\% & 8,5 & 56,8\% & 1,26 & 44,0 \\
Tolima & 1,58 & 54,5\% & 8,2 & 55,0\% & 0,58 & 44,3 \\
Caquetá & 1,57 & 54,1\% & 7,8 & 52,3\% & 0,15 & 46,3 \\
Boyacá & 1,56 & 54,0\% & 9,0 & 60,1\% & 0,52 & 40,1 \\
Norte de Santander & 1,50 & 52,0\% & 7,6 & 51,0\% & 0,67 & 45,6 \\
Huila & 1,40 & 48,5\% & 7,4 & 49,3\% & 0,47 & 44,0 \\
Cesar & 1,38 & 47,7\% & 7,1 & 47,8\% & 0,52 & 44,6 \\
Bolívar & 1,36 & 47,1\% & 7,2 & 48,2\% & 0,89 & 43,7 \\
Magdalena & 1,23 & 42,4\% & 6,5 & 43,4\% & 0,59 & 43,7 \\
Chocó & 1,22 & 42,1\% & 6,4 & 42,6\% & 0,14 & 44,2 \\
Cauca & 1,15 & 39,8\% & 6,9 & 45,9\% & 0,63 & 38,8 \\
Nariño & 1,10 & 38,1\% & 6,8 & 45,9\% & 0,85 & 37,2 \\
Córdoba & 1,08 & 37,4\% & 6,3 & 42,0\% & 0,79 & 39,9 \\
La Guajira & 1,02 & 35,1\% & 5,4 & 35,9\% & 0,35 & 43,7 \\
Sucre & 1,01 & 34,8\% & 6,0 & 40,2\% & 0,40 & 38,7 \\
\bottomrule
\end{tabular}
\caption*{\footnotesize Nota: remuneración por trabajador en millones de pesos al mes entre ocupados con ingreso horario positivo y horas válidas; remuneración por hora en miles de pesos para ese mismo universo. Ocupados en millones. Horas/trab. corresponde a horas semanales promedio por ocupado. Las columnas relativas comparan cada indicador con el departamento líder correspondiente (=100\%). Fuente: cálculos propios con GEIH.}
\end{table}